\chapter{Resultados}
\label{cap:resultados}

Este capítulo consolida os resultados quantitativos de validação
apresentados ao longo dos capítulos anteriores, organizados por nível de
integração: (i)~validação isolada da Fase~1 (propriedade de exatidão);
(ii)~validação campo a campo do \texttt{init.nc} completo contra um
arquivo real de produção; (iii)~validação do \texttt{lbc.*.nc}, por
auto-consistência e por comparação com a referência real; e
(iv)~validação funcional, comparando a previsão de 24 horas obtida a
partir dos arquivos desta rota com a previsão de referência gerada pela
rota original. Todos os números desta seção referem-se ao caso de estudo
usado consistentemente em todo o desenvolvimento: malha regional
\texttt{SouthAmerica} (recorte de $x1.163842$, 17064 células, 51648
arestas, $\sim\SI{60}{\kilo\metre}$ de resolução nominal), rodada iniciada
em \texttt{2026-01-01\_00:00:00}.

\section{Validação da interpolação horizontal (Fase 1)}

Conforme discutido na \cref{sec:validacao_fase1}, a interpolação de uma
malha de teste (2562 células) para os próprios centros de célula dessa
mesma malha (\emph{auto-\emph{scatter}}) produziu erro identicamente nulo
em 100\% das células, tanto para os valores de campo quanto para as
coordenadas de saída -- confirmando a implementação correta da
propriedade de exatidão nos vértices da interpolação baricêntrica
(\cref{sec:baricentrica}).

\section{Validação do \texttt{init.nc} completo}

\begin{table}[H]
\centering
\caption{Erro absoluto campo a campo entre o \texttt{init.nc} gerado por
  esta rota e o \texttt{init.nc} real de produção (17064 células, malha
  \texttt{SouthAmerica}, mesmo tempo inicial). Compilado e executado no
  build/ambiente real de produção (cluster Jaci), não apenas em ambiente
  de desenvolvimento local.}
\label{tab:validacao_init}
\begin{tabular}{@{}lrrr@{}}
\toprule
Campo & Faixa real (mín/máx) & Erro médio & Erro máximo \\
\midrule
$\rho$ [\si{\kilo\gram\per\meter\cubed}] & $0{,}018$ / $1{,}265$ & $1{,}5\times10^{-4}$ & $1{,}84\times10^{-2}$ \\
$\theta$ [\si{\kelvin}] & $273{,}1$ / $825{,}0$ & $0{,}114$ & $3{,}56$ \\
$q_v$ [\si{\kilo\gram\per\kilo\gram}] & $0$ / $0{,}0192$ & $5{,}1\times10^{-5}$ & $2{,}52\times10^{-3}$ \\
$u$ [\si{\meter\per\second}] & $-76$ / $77$ & $8{,}5\times10^{-2}$ & $4{,}47$ \\
$w$ [\si{\meter\per\second}] & $-1{,}52$ / $1{,}08$ & $4{,}0\times10^{-3}$ & $1{,}63$ \\
\texttt{surface\_pressure} [\si{\pascal}] & $5{,}76\times10^4$ / $1{,}03\times10^5$ & $5{,}9$ & $1792$ \\
\texttt{precipw} [\si{\milli\meter}] & $0{,}92$ / $62{,}1$ & $0{,}31$ & $4{,}46$ \\
\texttt{skintemp} [\si{\kelvin}] & $273{,}0$ / $307{,}5$ & $0{,}21$ & $7{,}07$ \\
\texttt{tmn} [\si{\kelvin}] & $266{,}2$ / $302{,}1$ & $0{,}15$ & $6{,}95$ \\
\texttt{vegfra} [\si{\percent}] & $0$ / $90{,}6$ & $7{,}3\times10^{-7}$ & $1{,}1\times10^{-5}$ \\
\bottomrule
\end{tabular}
\end{table}

Campos exatos ou essencialmente exatos (erro nulo ou de arredondamento de
armazenamento, $<10^{-6}$): \texttt{xland}, \texttt{xice},
\texttt{seaice}, \texttt{sh2o}, \texttt{dz}, \texttt{u\_init},
\texttt{v\_init}, \texttt{qv\_init}, \texttt{t\_init}, \texttt{qc},
\texttt{qr}, \texttt{sfc\_albbck}, \texttt{dzs}, \texttt{zs}. O arquivo
final contém a totalidade das 135 variáveis do \texttt{init.nc} real
(mais uma variável diagnóstica auxiliar, \texttt{psfc}, sem
correspondente no arquivo real e sem uso pelo núcleo do modelo). Os
resultados na \cref{tab:validacao_init} foram obtidos compilando e
executando o pipeline diretamente no ambiente de produção (compilador
GNU~13.3.0, bibliotecas NetCDF/PIO do cluster), reproduzindo,
\textbf{dígito a dígito}, os mesmos valores já obtidos em ambiente de
desenvolvimento local -- uma confirmação adicional de portabilidade e
robustez numérica da implementação frente a diferenças de compilador e
biblioteca.

\section{Validação do \texttt{lbc.*.nc}}

A auto-consistência entre o \texttt{init.nc} e um \texttt{lbc.*.nc}
gerado para o \emph{mesmo} tempo inicial (mesma malha vertical, mesmo
\textit{first-guess}) foi verificada como \textbf{exata}: diferença
identicamente nula em $\rho$, $\theta$, $q_v$, $u$ e $w$ -- o resultado
esperado, já que ambos os programas executam exatamente o mesmo cálculo
de Fases~3--4 para esse tempo.

\begin{table}[H]
\centering
\caption{Erro absoluto médio/máximo entre o \texttt{lbc.*.nc} gerado por
  esta rota e os arquivos reais de produção, nos cinco tempos de
  fronteira do experimento (intervalo de 6~h).}
\label{tab:validacao_lbc}
\begin{tabular}{@{}lrrrr@{}}
\toprule
Tempo & $u$ [\si{\meter\per\second}] & $\rho$ & $\theta$ [\si{\kelvin}] & $w$ [\si{\meter\per\second}] \\
\midrule
$t_0$ (00h)  & $0{,}086$ / $4{,}47$  & $0{,}0056$ / $0{,}043$ & $0{,}114$ / $3{,}56$ & $0{,}004$ / $1{,}63$ \\
$t_0+6$h  & $0{,}062$ / $5{,}69$  & $0{,}0056$ / $0{,}044$ & $0{,}055$ / $2{,}86$ & $0{,}004$ / $1{,}59$ \\
$t_0+12$h & $0{,}055$ / $7{,}10$  & $0{,}0056$ / $0{,}041$ & $0{,}051$ / $3{,}15$ & $0{,}004$ / $1{,}60$ \\
$t_0+18$h & $0{,}052$ / $12{,}2$  & $0{,}0057$ / $0{,}042$ & $0{,}060$ / $4{,}65$ & $0{,}004$ / $1{,}45$ \\
$t_0+24$h & $0{,}055$ / $12{,}5$  & $0{,}0057$ / $0{,}042$ & $0{,}059$ / $2{,}91$ & $0{,}004$ / $1{,}46$ \\
\bottomrule
\end{tabular}
\end{table}

O erro médio permanece pequeno e estável ao longo dos cinco tempos; o
erro máximo pontual cresce moderadamente nos tempos mais tardios --
compatível com divergência acumulada de previsão real ao longo de até 24
horas somada ao resíduo de fronteira já conhecido (\cref{cap:discussao}),
não com instabilidade numérica. O único problema identificado nesta
comparação não é desta implementação: o campo \texttt{lbc\_qv} do
arquivo de referência real é \textbf{constante em todos os 55 níveis
verticais por célula} -- fisicamente implausível para um campo de
umidade específica, e inconsistente com o comportamento correto (variável
por nível) dos demais campos do mesmo arquivo --, sugerindo um problema
já existente na cadeia de geração de referência (rota via WPS), não
atribuível a esta rota.

\section{Validação funcional: execução do modelo e comparação de previsão}

A previsão de 24 horas gerada a partir dos arquivos desta rota
(\cref{cap:fase7}) foi comparada, campo a campo e por tempo, com a
previsão de referência (mesma malha, mesmo período, gerada pela rota
original via WPS).

\begin{table}[H]
\centering
\caption{Erro absoluto médio/máximo entre a previsão gerada a partir dos
  arquivos desta rota e a previsão de referência, em três tempos de
  saída ao longo da integração de 24~h.}
\label{tab:validacao_previsao}
\begin{tabular}{@{}lrrrrr@{}}
\toprule
Tempo de previsão & $\theta$ [\si{\kelvin}] & $\rho$ & $q_v$ [\si{\kilo\gram\per\kilo\gram}] & $u$ [\si{\meter\per\second}] & $w$ [\si{\meter\per\second}] \\
\midrule
$+0$h  & $0{,}11$ / $3{,}6$ & $1{,}5\times10^{-4}$ / $1{,}8\times10^{-2}$ & $5{,}1\times10^{-5}$ / $2{,}5\times10^{-3}$ & $0{,}08$ / $4{,}5$ & $0{,}004$ / $1{,}6$ \\
$+12$h & $0{,}27$ / $6{,}4$ & $9{,}2\times10^{-4}$ / $2{,}9\times10^{-2}$ & $4{,}7\times10^{-4}$ / $1{,}8\times10^{-2}$ & $0{,}44$ / $9{,}2$ & $0{,}007$ / $0{,}33$ \\
$+24$h & $0{,}41$ / $9{,}8$ & $1{,}1\times10^{-3}$ / $3{,}0\times10^{-2}$ & $6{,}4\times10^{-4}$ / $1{,}8\times10^{-2}$ & $0{,}58$ / $12{,}0$ & $0{,}010$ / $0{,}65$ \\
\bottomrule
\end{tabular}
\end{table}

O padrão observado -- erro pequeno na inicialização, crescendo de forma
suave e limitada ao longo da integração, sem nenhum salto abrupto ou
sinal de instabilidade (ausência total de valores \texttt{NaN} em
qualquer campo, em qualquer tempo) -- é o esperado para duas trajetórias
vizinhas de um sistema dinâmico não linear sensível a condições iniciais
(divergência tipo Lyapunov): as duas rotas partem de condições iniciais
ligeiramente diferentes (pelo resíduo residual documentado nos
Capítulos~\ref{cap:fase2} e~\ref{cap:fase6}) e evoluem, sob a mesma
física e a mesma dinâmica, para estados progressivamente mais distantes
-- não uma divergência causada por um erro estrutural na implementação,
que se manifestaria como blowup numérico, quebra de conservação, ou
surgimento de valores não físicos, nenhum dos quais foi observado. O
próprio log de execução do modelo confirma a sanidade física ao longo de
toda a integração: os valores globais de mínimo/máximo de $w$ e $u$,
impressos a cada passo de tempo, permanecem dentro de faixas
meteorologicamente plausíveis (por exemplo, $w \in [-0{,}61,
+2{,}22]\,\si{\meter\per\second}$, $u \in [\pm 65, \pm 77]\,\si{\meter\per\second}$)
durante toda a janela de 24 horas.
